% Peyk - Model maliyeti ve acik agirlikli model alternatifleri
% Tablolar docs/cost/model.py tarafindan uretilir:  python3 model.py --latex
% Derleme:  pdflatex report.tex  (iki kez, icindekiler icin)
\documentclass[11pt,a4paper]{article}

\usepackage[utf8]{inputenc}
\usepackage[T1]{fontenc}
\usepackage[shorthands=off,turkish]{babel}
\usepackage{lmodern}
\usepackage{geometry}
\geometry{a4paper,margin=2.4cm,top=2.6cm,bottom=2.6cm}
\usepackage{booktabs}
\usepackage{array}
\usepackage{longtable}
\usepackage{xcolor}
\usepackage{enumitem}
\usepackage{titlesec}
\usepackage{fancyhdr}
\usepackage{caption}
\usepackage[hidelinks]{hyperref}
\hypersetup{pdftitle={Peyk --- Model maliyeti ve acik agirlikli model alternatifleri},
            pdfsubject={Teknik degerlendirme raporu},
            pdfkeywords={unit economics, Bedrock, open-weight models, triage}}

\definecolor{ink}{HTML}{16332C}
\definecolor{accent}{HTML}{8C5A0B}
\definecolor{soft}{HTML}{F1ECDF}
\definecolor{rule}{HTML}{9DB5AA}

\captionsetup{font=small,labelfont={bf,color=ink},skip=6pt}
\titleformat{\section}{\Large\bfseries\color{ink}}{\thesection.}{0.6em}{}
\titleformat{\subsection}{\large\bfseries\color{ink}}{\thesubsection}{0.6em}{}
\titlespacing*{\section}{0pt}{18pt}{8pt}

\pagestyle{fancy}\fancyhf{}
\renewcommand{\headrulewidth}{0.4pt}
\fancyhead[L]{\small\color{ink}Peyk --- Model maliyeti ve açık model alternatifleri}
\fancyhead[R]{\small\thepage}

\newcommand{\code}[1]{\texttt{\small #1}}
\newsavebox{\kbox}
\newenvironment{keybox}
  {\begin{center}\begin{lrbox}{\kbox}\begin{minipage}{0.90\textwidth}\color{ink}\small}
  {\end{minipage}\end{lrbox}\setlength{\fboxsep}{11pt}\colorbox{soft}{\usebox{\kbox}}\end{center}}

\begin{document}

\begin{titlepage}
\centering
\vspace*{3cm}
{\Huge\bfseries\color{ink} Peyk\par}
\vspace{0.6cm}
{\LARGE\color{ink} Model maliyeti ve açık ağırlıklı\\ model alternatifleri\par}
\vspace{0.8cm}
{\large\color{accent} Teknik değerlendirme raporu\par}
\vspace{2.2cm}
\begin{minipage}{0.8\textwidth}\small\color{ink}
\textbf{Kapsam.} Peyk'in kullanıcı başına model maliyetinin ölçülmesi, Claude Haiku~4.5
yerine açık ağırlıklı (open-weight) modellere geçiş senaryolarının değerlendirilmesi ve
maliyet--kalite dengesinin nicel olarak ortaya konması.
\vspace{0.5cm}

\textbf{Temel bulgu.} Açık modele geçiş, maliyet sıralamasında yanlış adımdır. Mevcut
kodda yapılacak ve hiçbir kalite riski taşımayan üç değişiklik aylık \$12{,}06 kazandırır;
triage'ı açık bir modele taşımak bunun üzerine yalnızca \$3{,}55 ekler ve ürünün tek
teknik iddiasını riske atar.
\end{minipage}
\vfill
{\small\color{ink} 15 Eylül 2026 \quad\textbullet\quad Ölçümler \code{akcware/peyk} deposundan\par}
\end{titlepage}

\tableofcontents
\newpage

% ==========================================================================
\section{Yönetici özeti}

Peyk'in bugünkü kullanıcı başına model maliyeti, orta yoğunlukta bir kullanıcı için
(günde 50 mail, 8 sohbet turu) \textbf{aylık \$19{,}30}'dur. Bu rakam \$25'lik bir abonelik
fiyatında \%23 brüt marj bırakır ki, Composio ve altyapı eklendiğinde sıfıra iner.

Sorunun kaynağı model \emph{seçimi} değil, model \emph{kullanımıdır}. Maliyetin
\%73'ü sohbet akışından gelir ve bu akış, her çağrıda 12.401 token'lık bir girdiyi
önbelleksiz olarak yeniden göndermektedir.

\begin{keybox}
\textbf{Bu raporun ana sonucu:} Takas edilmesi \emph{kolay} olan iş yükü (triage) ucuz
olandır; pahalı olan iş yükü (sohbet) ise takas edilmesi \emph{riskli} olandır. Açık
model geçişi, optimizasyonlardan önce yapılırsa üçte bir kadar kazanç için ürünün
güvenilirliği riske atılmış olur.
\end{keybox}

Önerilen sıra:

\begin{enumerate}[leftmargin=1.4em,itemsep=3pt]
  \item \textbf{Ölçüm altyapısı} (1 hafta). Her model çağrısının \code{usage} alanı
        veritabanına yazılmadan hiçbir karar veriye dayanmıyor.
  \item \textbf{Risksiz optimizasyonlar} (1--2 hafta). Prompt caching, bağlam küçültme,
        araç setinin bölünmesi, Sonnet~5'e geçiş. \$19{,}30 $\rightarrow$ \$7{,}24.
  \item \textbf{Değerlendirme setinin büyütülmesi} (2--3 hafta). Mevcut 60 mail'lik eval,
        model karşılaştırması için istatistiksel olarak yetersizdir.
  \item \textbf{Gölge modu ile aday ölçümü} (3--4 hafta). Ancak bundan sonra açık model
        kararı verilebilir; beklenen ek kazanç \$3{,}55.
\end{enumerate}

% ==========================================================================
\section{Yöntem ve kaynak güvenilirliği}

Bu raporun iki tür girdisi vardır ve ikisi çok farklı güvenilirlik taşır.

\subsection{Ölçülen veriler (yüksek güven)}

Token sayıları 15 Eylül 2026 tarihinde \code{akcware/peyk} deposundan doğrudan
çıkarılmıştır. Kaynaklar Ek~B'de dosya ve satır numarasıyla listelenmiştir. Sayımlar
karakter uzunluğunun 4'e bölünmesiyle yapılmıştır; gerçek tokenizer çıktısına göre
\textbf{$\pm$\%30} sapma payı vardır.

\subsection{Varsayılan veriler (düşük güven, doğrulanmalı)}

\begin{keybox}
\textbf{Uyarı.} Bu raporun hazırlandığı ortamın ağ politikası fiyat kaynaklarının
tamamını engellemiştir: \code{aws.amazon.com/bedrock/pricing},
\code{docs.aws.amazon.com}, \code{openrouter.ai}, \code{artificialanalysis.ai},
\code{huggingface.co} ve \code{together.ai} çağrıları 403 ile reddedilmiştir.
Hiçbir fiyat canlı kaynaktan doğrulanamamıştır.
\end{keybox}

Bu kısıt nedeniyle rapor iki savunma mekanizması kullanır:

\begin{itemize}[leftmargin=1.4em,itemsep=3pt]
  \item \textbf{Anthropic fiyatları} liste fiyatlarıdır (Haiku~4.5 \$1/\$5,
        Sonnet~4.6 \$3/\$15, Sonnet~5 \$2/\$10 --- milyon token başına giriş/çıkış).
        Bedrock ayrı fiyatlandırılan bir iş ortağı platformudur; Claude modellerinde
        tarihsel olarak liste fiyatıyla örtüşmüştür ancak \textbf{bu doğrulanmalıdır}.
  \item \textbf{Açık model fiyatları tedarikçi iddiası olarak değil, fiyat bandı olarak}
        modellenmiştir (Bant A--D, \$0{,}05--\$0{,}70). Böylece sonuçlar belirli bir
        satıcının güncel fiyatına bağlı kalmaz. Hangi modelin hangi banda düştüğü
        Bölüm~6'da ayrıca tartışılmıştır ve doğrulanması gereken tek şey budur.
\end{itemize}

Tüm hesaplar \code{docs/cost/model.py} dosyasındadır. Fiyatlar tek bir sözlükte
parametredir; güncellenip çalıştırıldığında bu rapordaki her tablo yeniden üretilir.

% ==========================================================================
\section{Ölçülen token profili}

\begin{table}[h]
\centering\small
\begin{tabular}{@{}lrrl@{}}
\toprule
\textbf{Çağrı} & \textbf{Girdi} & \textbf{Çıktı} & \textbf{Sıklık} \\
\midrule
Triage (Haiku 4.5) & 1.740 & 150 & her dış olay \\
Ack / refleks (Haiku 4.5) & 1.330 & 60 & her sohbet turu \\
Sohbet (Sonnet 4.6) & 12.401 & 250 & her sohbet turu $\times$ 1,4 \\
\bottomrule
\end{tabular}

\caption{Çağrı başına token profili. Ack ve sohbet her kullanıcı mesajında birlikte çalışır.}
\end{table}

Buradan çıkan birim maliyetler: \textbf{bir mail triage'ı \$0{,}0025},
\textbf{bir sohbet turu \$0{,}0590}. Depodaki \code{docs/DEPLOY.md} dosyası triage için
``\$0{,}001 per mail'' demektedir; gerçek rakam bunun 2{,}5 katıdır. Fark, mailin
kendisinden değil, her çağrıda yeniden gönderilen 1.140 token'lık sistem promptundan
kaynaklanır.

\subsection{Görev bazında maliyet payı}

\begin{table}[h]
\centering\small
\begin{tabular}{@{}lrr@{}}
\toprule
\textbf{Görev} & \textbf{\$/ay} & \textbf{Pay} \\
\midrule
Triage & 3,73 & 19,4\% \\
Sohbet & 14,15 & 73,3\% \\
Sabah brief & 1,42 & 7,3\% \\
\bottomrule
\end{tabular}

\caption{Orta profil (50 mail, 8 sohbet/gün), bugünkü yapılandırma.}
\end{table}

Bu tablo raporun geri kalanını belirler: \textbf{triage toplam maliyetin yalnızca
\%19{,}4'üdür.} Triage'ı tamamen bedava bir modele taşımak bile, en iyi ihtimalle aylık
\$3{,}73 kazandırır.

% ==========================================================================
\section{Maliyet nerede yoğunlaşıyor}

Bir sohbet turunun 12.401 token'lık girdisinin dağılımı:

\begin{table}[h]
\centering\small
\begin{tabular}{@{}lrr@{}}
\toprule
\textbf{Bileşen} & \textbf{Token} & \textbf{Pay} \\
\midrule
\code{recent\_observations} (limit 50) & 5.000 & 40,3\% \\
23 adet \code{@tool} şeması & 3.350 & 27,0\% \\
\code{CHAT\_SYSTEM\_PROMPT} & 3.126 & 25,2\% \\
Geçmiş 10 tur & 625 & 5,0\% \\
Hafıza + kullanıcı + iskele & 300 & 2,4\% \\
\bottomrule
\end{tabular}

\caption{Sohbet girdisinin bileşenleri (\code{workers/chat.py}).}
\end{table}

Üç yapısal sorun görünüyor:

\begin{enumerate}[leftmargin=1.4em,itemsep=4pt]
  \item \textbf{Bağlam peşin yükleniyor.} \code{RECENT\_LIMIT = 50}, 48 saatlik 50 gözlemi
        her sohbete önden koyuyor --- girdinin \%40'ı. Oysa ajanda zaten \code{need\_more}
        aracı var; veriyi isteyebiliyor. Bu, tek satırlık bir değişiklikle düzelir ve
        önbelleklemeye hiç bağlı değildir.
  \item \textbf{23 araç şeması her çağrıda gidiyor.} Takvimle ilgili altı araç
        (\code{find\_events}, \code{find\_free\_time}, \code{check\_time},
        \code{create\_event}, \code{update\_event}, \code{cancel\_event}) yalnızca
        mesajda bir zaman geçtiğinde gereklidir.
  \item \textbf{Hiçbir yerde prompt caching yok.} Sistem promptu ve araç şemaları birlikte
        6.476 token'lık, çağrıdan çağrıya \emph{bayt bayt aynı} bir önek oluşturuyor.
\end{enumerate}

\subsection{Önbelleklemede sessiz bir tuzak}

Prompt caching'in modele göre değişen bir minimum önek uzunluğu vardır; bu eşiğin altında
önbellek \emph{sessizce} çalışmaz --- hata verilmez, yalnızca
\code{cache\_creation\_input\_tokens} sıfır döner.

\begin{table}[h]
\centering\small
\begin{tabular}{@{}llrl@{}}
\toprule
\textbf{Çağrı} & \textbf{Model} & \textbf{Önek} & \textbf{Minimum / sonuç} \\
\midrule
Sohbet & Sonnet 4.6 & 6.476 tok & 1.024 --- çalışır \\
Triage & Haiku 4.5 & 1.340 tok & \textbf{4.096 --- sessizce çalışmaz} \\
\bottomrule
\end{tabular}
\caption{Haiku 4.5'in minimum önbelleklenebilir önek uzunluğu 4.096 token'dır.}
\end{table}

Sonuç: sohbet önbelleklenebilir, \textbf{triage önbelleklenemez}. Bu, triage'ın birim
maliyetini \$0{,}0025'te sabitler ve açık model arayışının asıl gerekçesini oluşturur.

İlginç bir ara sonuç: triage'ı önbellekli Sonnet~5 ile çalıştırmak \$0{,}0026'ya mal olur
--- yani \emph{aynı paraya belirgin biçimde daha güçlü bir model}. Ancak önbelleğin
5~dakikalık yaşam süresi, seyrek gelen maillerde önbelleğin çoğu zaman soğuk olması
anlamına gelir; bu durumda yazma cezası (1{,}25$\times$) devreye girer ve avantaj kaybolur.
Sohbet turları kümelendiği için orada aynı sorun yoktur.

% ==========================================================================
\section{Hangi iş yükü takas edilebilir?}

Her model çağrısının farklı bir gereksinim profili vardır. Takas edilebilirlik, fiyattan
önce bu profille belirlenir.

\begin{table}[h]
\centering\small
\renewcommand{\arraystretch}{1.25}
\begin{tabular}{@{}lccccl@{}}
\toprule
\textbf{Görev} & \textbf{Yapısal} & \textbf{Araç} & \textbf{Çok} & \textbf{Bağlam} & \textbf{Takas} \\
 & \textbf{çıktı} & \textbf{çağrısı} & \textbf{dilli} & & \textbf{edilebilirlik} \\
\midrule
Triage & evet & hayır & evet & 1{,}7K & \textbf{Yüksek} \\
Ack / refleks & evet & hayır & evet & 1{,}3K & \textbf{Çok yüksek} \\
Learn & evet & hayır & evet & $\sim$8K & Yüksek (seyrek) \\
Sohbet & hayır & \textbf{23 araç} & evet & 12{,}4K & \textbf{Düşük} \\
Gömme (Titan v2) & --- & --- & --- & --- & Zaten ucuz \\
\bottomrule
\end{tabular}
\caption{İş yükü profilleri. Sohbet dışındaki her şey tek atışlık yapısal çıktıdır.}
\end{table}

\textbf{Triage ve ack} tek atışlık sınıflandırma görevleridir: sabit şema, kısa çıktı,
araç yok. Açık ağırlıklı küçük modellerin en güçlü olduğu alan tam da budur.

\textbf{Sohbet} ise tersidir. 23 aracın arasından doğru olanı seçmek, çok turlu araç
çağrısı zincirini bozmadan yürütmek ve bunu kullanıcının dilinde yapmak, küçük açık
modellerin sistematik olarak zorlandığı bir iştir. Üstelik burada bir hata sessiz değildir:
yanlış araç çağrısı yanlış bir taslak maile veya yanlış bir takvim değişikliğine dönüşür.
Onay akışı (\code{approval.py}) bunu büyük ölçüde yakalar, ama kullanıcı deneyimi bozulur.

\begin{keybox}
Maliyetin \%73'ü sohbette, takas kolaylığının tamamı triagededir. Bu raporun merkezi
gerilimi budur: \textbf{ucuzlatması kolay olan yer az para tutuyor.}
\end{keybox}

% ==========================================================================
\section{Açık ağırlıklı model manzarası}

\subsection{Bedrock içinde kalan seçenekler}

Bedrock, Claude dışında birden fazla açık ağırlıklı model ailesi barındırır: Meta Llama
(3.x ve 4 serisi), Mistral, DeepSeek, Qwen ve \code{gpt-oss} aileleri; ayrıca açık
ağırlıklı olmayan ama çok ucuz olan Amazon Nova serisi. Bunların Peyk açısından belirleyici
avantajı fiyat değil, \textbf{tedarikçi değişmemesidir}: aynı IAM anahtarı, aynı bölge,
aynı veri işleme sözleşmesi, GDPR alt işleyici listesinde yeni bir satır yok.

\subsection{Bedrock dışı seçenekler ve gizli maliyeti}

Groq, Together, Fireworks, DeepInfra ve OpenRouter gibi sağlayıcılar aynı açık modelleri
genellikle daha ucuza sunar. Ancak Peyk için bu \textbf{yalnızca bir fiyat kararı değildir}:

\begin{itemize}[leftmargin=1.4em,itemsep=3pt]
  \item Triage girdisi kullanıcının \emph{mail gövdesini} içerir
        (\code{mappings.py}, \code{message\_text} alanı, 1.200 karakteri kapsar).
        Yeni bir sağlayıcı, kullanıcının mail içeriğine erişen yeni bir alt işleyicidir.
  \item Bu, GDPR alt işleyici listesini, kullanıcıya açıklanan gizlilik politikasını ve
        Google'ın kısıtlı kapsam (restricted scope) değerlendirmesinde beyan edilen veri
        akışını doğrudan etkiler.
  \item Aylık \$3{,}55'lik bir tasarruf için yeni bir DPA müzakeresi, yeni bir
        güvenlik değerlendirmesi ve gizlilik politikası revizyonu gerekir.
\end{itemize}

\begin{keybox}
\textbf{Öneri:} Açık model denemeleri Bedrock içinde kalsın. Fiyat avantajının bir kısmı
kaybedilir, ancak hukuki ve operasyonel maliyet sıfırdır. Bedrock dışına çıkmak ancak
kullanıcı sayısı bu ek yükü haklı çıkardığında --- kabaca birkaç bin kullanıcıda ---
gündeme gelmelidir.
\end{keybox}

\subsection{Fiyat bantları}

Belirli tedarikçi fiyatları doğrulanamadığı için modelleme dört bant üzerinden yapılmıştır.
Her bandın karşılığı olarak verilen model sınıfları \textbf{örnektir, doğrulanmalıdır}.

\begin{table}[h]
\centering\small
\begin{tabular}{@{}llll@{}}
\toprule
\textbf{Bant} & \textbf{Giriş/Çıkış (\$/MTok)} & \textbf{Tipik model sınıfı} & \textbf{Triage için} \\
\midrule
A & 0{,}05 / 0{,}20 & Nova Micro sınıfı & Kalite riski yüksek \\
B & 0{,}15 / 0{,}30 & 3B--20B açık modeller & Aday \\
C & 0{,}35 / 0{,}60 & 7--8B açık modeller & \textbf{En dengeli aday} \\
D & 0{,}70 / 0{,}90 & 70B açık modeller & Güvenli ama az kazanç \\
\midrule
--- & 1{,}00 / 5{,}00 & Claude Haiku 4.5 (mevcut) & Referans \\
\bottomrule
\end{tabular}
\caption{Fiyat bantları. Model sınıfı eşlemeleri doğrulanmayı bekleyen varsayımlardır.}
\end{table}

% ==========================================================================
\section{Senaryo matrisi}

Aşağıdaki tablo kümülatiftir: her satır bir öncekinin üzerine ekler.

\begin{table}[h]
\centering\small
\renewcommand{\arraystretch}{1.2}
\begin{tabular}{@{}llrrr@{}}
\toprule
& \textbf{Senaryo} & \textbf{Hafif} & \textbf{Orta} & \textbf{Yoğun} \\
& & \multicolumn{3}{c}{\footnotesize \$/kullanıcı/ay} \\
\midrule
\textbf{S0} & Bugünkü hali & 8,22 & 19,30 & 45,76 \\
\textbf{S1} & Prompt caching açık & 5,43 & 12,84 & 30,48 \\
\textbf{S2} & S1 + gözlem 50 $\rightarrow$ 15 & 3,75 & 8,96 & 21,31 \\
\textbf{S3} & S2 + araç seti bölünmüş & 3,67 & 8,78 & 20,88 \\
\textbf{S4} & S3 + Sonnet 5 & 3,01 & 7,24 & 17,25 \\
\textbf{S5} & S4 + triage Bant C (\$0,35) & 1,68 & 3,92 & 9,28 \\
\textbf{S6} & S4 + triage Bant B (\$0,15) & 1,59 & 3,69 & 8,73 \\
\textbf{S7} & S6 + ack de açık model & 1,43 & 3,32 & 7,85 \\
\textbf{S8} & S7 + sohbet de Bant D (riskli) & 0,46 & 1,08 & 2,56 \\
\bottomrule
\end{tabular}

\caption{Kümülatif senaryo matrisi. Hafif = 20 mail + 3 sohbet/gün; Orta = 50 + 8;
Yoğun = 120 + 20.}
\end{table}

Okunması gereken üç eşik:

\begin{description}[leftmargin=0pt,style=unboxed,itemsep=5pt]
  \item[S0 $\rightarrow$ S4: \$19{,}30 $\rightarrow$ \$7{,}24.] Tamamı mevcut kodda
        yapılacak değişikliklerdir. Model değişmez, kalite riski yoktur, tek fark
        Sonnet~4.6 yerine hem daha yeni hem daha ucuz olan Sonnet~5'tir.
        \textbf{Kazanç: \$12{,}06.}
  \item[S4 $\rightarrow$ S6: \$7{,}24 $\rightarrow$ \$3{,}69.] Triage'ın açık modele
        taşınması. \textbf{Kazanç: \$3{,}55}, ve bu kazanç için ürünün tek teknik iddiası
        ölçüm altına alınmak zorundadır.
  \item[S7 $\rightarrow$ S8: \$3{,}32 $\rightarrow$ \$1{,}08.] Sohbetin de açık modele
        taşınması. Mutlak kazanç büyüktür (\$2{,}24) çünkü maliyetin \%73'ü oradadır ---
        ama 23 araçlı çok turlu akışta bu, ürünün davranışını baştan doğrulamayı gerektirir.
        Bu rapor bunu önermez; yalnızca üst sınırı göstermek için tabloda tutulmuştur.
\end{description}

% ==========================================================================
\section{Triage fiyat duyarlılığı}

\begin{table}[h]
\centering\small
\begin{tabular}{@{}lrrrr@{}}
\toprule
\textbf{Model bandı} & \textbf{\$/MTok} & \textbf{Birim} & \textbf{\$/ay} & \textbf{Tasarruf} \\
\midrule
Claude Haiku 4.5 & 1,00 / 5,00 & 0,00249 & 3,73 & 0,00 \\
Bant D: $\sim$\$0,70 & 0,70 / 0,90 & 0,00135 & 2,03 & 1,71 \\
Bant C: $\sim$\$0,35 & 0,35 / 0,60 & 0,00070 & 1,05 & 2,69 \\
Bant B: $\sim$\$0,15 & 0,15 / 0,30 & 0,00031 & 0,46 & 3,28 \\
Bant A: $\sim$\$0,05 & 0,05 / 0,20 & 0,00012 & 0,18 & 3,56 \\
\bottomrule
\end{tabular}

\caption{Orta profil (50 mail/gün). ``Tasarruf'' Haiku 4.5'e göredir.}
\end{table}

Bu tablonun en önemli özelliği \textbf{azalan getiridir}. Bant~D'den Bant~A'ya inmek ---
yani 14 kat daha ucuz bir modele geçmek --- yalnızca \$1{,}85 ek kazandırır. Maliyetin
büyük kısmı zaten sabit sistem promptundadır ve o token'lar hangi modelde olursa olsun
gönderilir.

\begin{keybox}
``En ucuz modeli bulma'' yarışının anlamı yoktur. Bant~C'nin (\$0{,}35) altına inmenin
getirisi \$1{,}05/ay'dır; buna karşılık model küçüldükçe kalite riski hızla artar.
\textbf{Aday aralığı Bant C--D'dir.}
\end{keybox}

% ==========================================================================
\section{Kendi altyapında barındırma (self-hosting)}

Açık ağırlıklı modelin asıl cazibesi marjinal maliyetin sıfıra yaklaşmasıdır. Ancak sabit
maliyet GPU kirasıdır ve başabaş noktası şaşırtıcı biçimde yüksektir.

\begin{table}[h]
\centering\small
\begin{tabular}{@{}lrrlrr@{}}
\toprule
\textbf{Örnek tipi} & \textbf{\$/sa} & \textbf{\$/ay} & \textbf{Sığan model} & \textbf{Başabaş} & \textbf{Kullanıcı} \\
\midrule
g5.xlarge (A10G 24GB) & 1,006 & 724 & 7-8B, 4-bit & 290,892 & 194 \\
g6.2xlarge (L4 24GB) & 0,978 & 704 & 7-8B & 282,795 & 189 \\
g5.12xlarge (4xA10G) & 5,672 & 4.084 & 70B, 4-bit & 1,640,096 & 1.093 \\
\bottomrule
\end{tabular}

\caption{Başabaş, Haiku 4.5 ile triage maliyetine göre. On-demand fiyatlar varsayımdır.}
\end{table}

Yani \textbf{yaklaşık 190 aktif kullanıcının altında kendi GPU'nu çalıştırmak
doğrudan zarardır} --- ve bu hesap yalnızca GPU kirasıdır. Dahil olmayanlar: model
sunucusunun (vLLM vb.) işletimi, ölçekleme, yedeklilik, gece nöbeti, model güncellemeleri
ve GPU kotası temini.

Ayrıca Peyk'in mevcut mimarisi tek Fly.io makinesine kilitlidir (\code{fly.toml},
\code{auto\_stop\_machines = off}) çünkü Telegram \code{getUpdates} ve Composio websocket
tek tüketicidir. Kendi model sunucusunu eklemek, henüz yatay ölçeklenemeyen bir sisteme
ikinci bir durum bilgili bileşen eklemek anlamına gelir.

\begin{keybox}
Self-hosting bu aşamada gündem dışıdır. 200+ ödeyen kullanıcıya ulaşıldığında, ve ancak
Bölüm~2'deki webhook geçişi tamamlandıktan sonra yeniden değerlendirilmelidir.
\end{keybox}

% ==========================================================================
\section{Kalite riski: asimetrik maliyet}

Ucuz model kararı, tasarrufu kalite riskine karşı tartmadan verilemez. Peyk'te bu denge
özellikle asimetriktir çünkü ürünün tüm vaadi ``doğru olanı geçiriyorum'' iddiasıdır.

İki hata türü eşit değildir:

\begin{itemize}[leftmargin=1.4em,itemsep=3pt]
  \item \textbf{Yanlış pozitif} (önemsizi acil sanmak): kullanıcı bir gereksiz bildirim alır,
        \code{noise} basar, bütçe slotu geri verilir (\code{budget\_repo.py}). Maliyeti düşüktür.
  \item \textbf{Yanlış negatif} (gerçek acili kaçırmak): kullanıcı müşterisini, kirasını
        veya çocuğunun okulundan gelen çağrıyı kaçırır. Bu tek olay, ürüne olan güveni
        bitirir.
\end{itemize}

Aşağıdaki tablo, bir tasarrufun hangi ek kaçırma oranında sıfırlandığını gösterir.
Varsayımlar: abonelik \$25/ay, ortalama ömür 12 ay (LTV \$300), kaçırılan bir acil olayın
ayrılmaya yol açma olasılığı \%50.

\begin{table}[h]
\centering\small
\begin{tabular}{@{}lrr@{}}
\toprule
\textbf{Model bandı} & \textbf{Tasarruf \$/ay} & \textbf{Başabaş ek kaçırma} \\
\midrule
Bant D: $\sim$\$0,70 & 1,71 & 1,14\% \\
Bant C: $\sim$\$0,35 & 2,69 & 1,79\% \\
Bant B: $\sim$\$0,15 & 3,28 & 2,18\% \\
Bant A: $\sim$\$0,05 & 3,56 & 2,37\% \\
\bottomrule
\end{tabular}

\caption{``Başabaş ek kaçırma'', 100 kullanıcı-ayında kaç ek kaçırılmış acilin tasarrufu
sıfırladığını gösterir.}
\end{table}

Okuması şudur: Bant~B'ye geçmek, \textbf{100 kullanıcı-ayında 2{,}2 ek kaçırılmış acile
mal olursa} tasarruf tamamen silinir. Bu, çok dar bir hata bütçesidir --- ve tam olarak
ölçülmesi gereken şeydir.

% ==========================================================================
\section{Karar protokolü}

Peyk'in bu kararı ampirik olarak verebilecek altyapısı büyük ölçüde hazırdır.
\code{tests/phase1/}\allowbreak\code{test\_triage\_eval.py} 60 etiketli mail üzerinde çalışır; mevcut ölçüm
\%98 within-1 doğruluk ve 6/6 urgency-5 recall'dur.

\subsection{Önce: değerlendirme seti büyütülmeli}

60 mail, model \emph{karşılaştırması} için yetersizdir. Kritik metrik olan urgency-5
recall yalnızca 6 örneğe dayanmaktadır; bir modelin 6/6, diğerinin 5/6 alması istatistiksel
olarak anlamlı bir fark değildir. Bölüm~10'daki hata bütçesini ölçebilmek için
\textbf{en az 300 etiketli mail ve bunların içinde en az 30 urgency-5 örneği} gerekir.

Bu setin kaynağı hazırdır: \code{sent\_notification.user\_feedback} alanında biriken
\code{useful}/\code{noise} geri bildirimleri ve gate'in sustuğu olaylar. Ayrıca bu veri
zamanla rakiplerin sahip olmadığı bir varlığa dönüşür --- küçük bir modelin bu veriyle
ince ayarlanması (fine-tuning), uzun vadede hem maliyet hem savunulabilirlik avantajıdır.

\subsection{Sonra: gölge modu}

Aday modeller üretimde şu şekilde risksiz ölçülebilir:

\begin{enumerate}[leftmargin=1.4em,itemsep=3pt]
  \item Her gerçek triage çağrısı \emph{ayrıca} aday modele de gönderilir.
  \item Kullanıcıya \textbf{yalnızca} Haiku'nun sonucu gider; aday modelin çıktısı sadece
        loglanır.
  \item İki model arasındaki farklar, özellikle Haiku'nun 4--5 verdiği ama adayın 1--3
        verdiği olaylar, elle incelenir.
\end{enumerate}

Gölge modunun maliyeti aday modelin fiyatı kadardır (Bant~C için aylık \$1{,}05) ve
kullanıcı için sıfır risk taşır. İki hafta çalıştırmak, bu raporun tüm varsayımlarını
gerçek veriyle değiştirir.

\subsection{Kabul kriterleri}

\begin{table}[h]
\centering\small
\begin{tabular}{@{}lll@{}}
\toprule
\textbf{Metrik} & \textbf{Eşik} & \textbf{Gerekçe} \\
\midrule
Urgency-5 recall & \textbf{\%100} & Kaçırılan acil ürünü bitirir \\
Within-1 doğruluk & $\geq$ \%95 & Mevcut: \%98 \\
Yapısal çıktı hata oranı & $<$ \%0{,}5 & Şema uyumsuzluğu = düşen olay \\
Türkçe/Almanca özet kalitesi & elle kontrol & Çok dilli özet küçük modellerde zayıflar \\
p95 gecikme & $<$ 3 sn & Gate zaten dakikalarca bekliyor, esnek \\
\bottomrule
\end{tabular}
\caption{Bir aday modelin üretime alınabilmesi için karşılaması gereken eşikler.}
\end{table}

% ==========================================================================
\section{Sonuç ve yol haritası}

\begin{table}[h]
\centering\small
\renewcommand{\arraystretch}{1.3}
\begin{tabular}{@{}llll@{}}
\toprule
\textbf{Faz} & \textbf{İş} & \textbf{Süre} & \textbf{Sonuç} \\
\midrule
0 & \code{usage} muhasebesi (\code{model\_call} tablosu) & 1 hafta & Varsayım $\rightarrow$ veri \\
1 & Caching + bağlam + araç bölme + Sonnet 5 & 1--2 hafta & \$19{,}30 $\rightarrow$ \$7{,}24 \\
2 & Eval seti 60 $\rightarrow$ 300+ mail & 2--3 hafta & Ölçülebilir karar \\
3 & Gölge modu, Bant C--D adayları & 3--4 hafta & Kanıt \\
4 & Karar & --- & Olası \$7{,}24 $\rightarrow$ \$3{,}69 \\
\bottomrule
\end{tabular}
\caption{Önerilen sıralama.}
\end{table}

Faz~1 ile Faz~4'ün karşılaştırması bu raporun özetidir: birincisi \$12{,}06 kazandırır ve
hiçbir şeyi riske atmaz; ikincisi \$3{,}55 kazandırır ve ürünün tek teknik iddiasını
ölçüm altına almayı zorunlu kılar. İkisi de yapılmalıdır --- ama bu sırayla.

\$25'lik fiyat noktasında brüt marjın seyri:

\begin{table}[h]
\centering\small
\begin{tabular}{@{}lrrr@{}}
\toprule
\textbf{Aşama} & \textbf{Hafif} & \textbf{Orta} & \textbf{Yoğun} \\
\midrule
Bugün (S0) & \%67 & \%23 & \textbf{--\%83} \\
Faz 1 sonrası (S4) & \%88 & \%71 & \%31 \\
Faz 4 sonrası (S6) & \%94 & \%85 & \%65 \\
\bottomrule
\end{tabular}
\caption{Brüt marj, \$25/ay abonelik. Composio ve altyapı hariç.}
\end{table}

Yoğun kullanıcı satırı ayrıca dikkat ister: bugün ürüne en çok değer veren kullanıcı en
çok para kaybettiren kullanıcıdır. Faz~1 bunu düzeltir, ancak uzun vadede bir adil
kullanım sınırı gerekecektir. Peyk'in ürün anlatısı buna elverişlidir --- knock bütçesinin
yanına bir sohbet bütçesi koymak, ürünün hikâyesiyle çelişmez.

\appendix
% ==========================================================================
\section{Modeli yeniden çalıştırma}

\begin{verbatim}
cd docs/cost
python3 model.py              # okunabilir tablolar
python3 model.py --latex      # LaTeX tablo gövdeleri
pdflatex report.tex           # (iki kez)
\end{verbatim}

Fiyatlar \code{model.py} içindeki \code{PRICES} sözlüğündedir. Bedrock'un gerçek
fiyatları doğrulandığında yalnızca o sözlük güncellenir; bu rapordaki her tablo yeniden
üretilir. Aynı şekilde \code{PROFILES} (kullanım yoğunluğu) ve
\code{CHURN\_GIVEN\_MISS} / \code{LTV} (risk varsayımları) da parametredir.

\section{Token sayımlarının kaynağı}

\begin{table}[h]
\centering\small
\begin{tabular}{@{}llr@{}}
\toprule
\textbf{Değer} & \textbf{Kaynak} & \textbf{Ölçüm} \\
\midrule
Triage sistem promptu & \code{agent/triage\_agent.py} & 4.151 karakter \\
Sohbet sistem promptu & \code{agent/chat\_agent.py} & 12.504 karakter \\
Ack sistem promptu & \code{agent/chat\_agent.py} & 3.121 karakter \\
23 araç şeması & \code{agent/chat\_agent.py} & 10.358 karakter \\
Gözlem limiti & \code{workers/chat.py} & \code{RECENT\_LIMIT = 50} \\
Geçmiş tur sayısı & \code{workers/chat.py} & \code{HISTORY\_TURNS = 10} \\
Bağlam penceresi & \code{workers/chat.py} & \code{RECENT\_HOURS = 48} \\
Mail gövdesi kırpma & \code{agent/triage\_agent.py} & \code{max\_text = 1200} \\
Tur sayısı & \code{workers/chat.py} & \code{MAX\_ROUNDS = 2}, \code{MAX\_DOC\_ROUNDS = 4} \\
\bottomrule
\end{tabular}
\caption{Ölçümlerin depo içindeki kaynakları (15 Eylül 2026).}
\end{table}

\vspace{1cm}
\begin{center}
\small\color{ink}
\textit{Bu rapordaki fiyat varsayımları canlı kaynaktan doğrulanmamıştır (bkz. Bölüm 2).\\
Herhangi bir tedarikçi kararından önce güncel fiyatlar teyit edilmelidir.}
\end{center}

\end{document}
